\DocumentMetadata{
  pdfversion=2.0,
  pdfstandard=ua-2,
  lang=en-US,
  tagging=on,
  tagging-setup={math/setup=mathml-SE,tabsorder=structure}
}
\documentclass[11pt,letterpaper]{article}
\usepackage[margin=0.68in,includefoot]{geometry}
\usepackage{newcomputermodern}
\usepackage{amsmath,amscd,mathtools}
\usepackage{tikz,adjustbox}
\providecommand{\square}{\mdlgwhtsquare}
\usepackage[hidelinks]{hyperref}
\hypersetup{
  pdftitle={Analysis First-Year Exam, January 2014, Part 1},
  pdfauthor={University of Florida Department of Mathematics},
  pdfsubject={Graduate mathematics examination},
  pdfdisplaydoctitle=true,
  bookmarksopen=true
}
\setcounter{secnumdepth}{0}
\setlength{\parindent}{0pt}
\setlength{\parskip}{0.28em}
\setlength{\emergencystretch}{3em}
\setlength{\leftmargini}{2.15em}
\setlength{\leftmarginii}{2.65em}
\setlength{\leftmarginiii}{3em}
\renewcommand{\labelenumi}{\textbf{\theenumi.}}
\pagestyle{empty}
\raggedbottom
\allowdisplaybreaks
\newenvironment{examproblems}[1]{%
  \begin{enumerate}%
  \setcounter{enumi}{#1}%
  \setlength{\topsep}{0.35em}%
  \setlength{\partopsep}{0pt}%
  \setlength{\itemsep}{0.48em}%
  \setlength{\parsep}{0.18em}%
}{\end{enumerate}}
\newenvironment{examsubparts}{%
  \begin{enumerate}%
  \setlength{\topsep}{0.18em}%
  \setlength{\partopsep}{0pt}%
  \setlength{\itemsep}{0.16em}%
  \setlength{\parsep}{0.1em}%
}{\end{enumerate}}
\newenvironment{examinstructions}{%
  \par\begingroup\itshape
}{%
  \par\endgroup\vspace{0.18em}
}
\makeatletter
\renewcommand\section{\@startsection{section}{1}{0pt}%
  {-1.25ex plus -.25ex minus -.1ex}%
  {0.55ex plus .1ex}%
  {\normalfont\large\bfseries}}
\renewcommand{\@maketitle}{%
  \newpage\null\vskip -1em%
  {\footnotesize\bfseries
    \href{https://www.ufl.edu/}{University of Florida}, \href{https://math.ufl.edu/}{Department of Mathematics}\par}%
  \vskip -0.5em%
  \tagstructbegin{tag=Title}%
    {\LARGE\bfseries\href{https://gma.math.ufl.edu/exams/analysis-first-year/}{Analysis First-Year Exam}, January 2014, Part 1\par}%
  \tagstructend%
  \vskip 0.25em%
  \tagmcbegin{artifact}%
    \hrule height 1pt%
  \tagmcend%
  \vskip 1em%
}
\makeatother
\title{Analysis First-Year Exam, January 2014, Part 1}
\author{}
\date{}
\begin{document}
\maketitle
\thispagestyle{empty}
\begin{examinstructions}
Do only 4 problems of the 6. Work must be presented in a neat and logical fashion to receive credit. Do not leave any gaps. State clearly any theorems used in proofs.
\end{examinstructions}
\begin{examproblems}{0}
\item
Let~\(F\) be a closed subset of the reals~\(\mathbb{R}\) with the usual topology. Show~\(F\) is a countable union of compact sets.
\item
Let~\(E\) be a nonempty subset of a metric space. State the definition of the distance \(\rho_E(x)\) of a point~\(x\) to~\(E\). Characterize \(\overline{E}\) in terms of \(\rho_E\), where \(\overline{E}\) is the closure of~\(E\).
\item
This problem has two parts.
\begin{examsubparts}
\item[\textnormal{(a)}]
Let~\(f\) be a continuous function from a compact metric space~\(X\) into a metric space~\(Y\). Prove \(f(X)\) is compact.
\item[\textnormal{(b)}]
Assume the setting in (a) with the further assumption that~\(Y\) is the reals with the usual topology. Show that~\(f\) assumes maximum and minimum values on~\(X\).
\end{examsubparts}
\item
Let~\(f\) be a continuous map of a metric space~\(X\) into a metric space~\(Y\).
\begin{examsubparts}
\item[\textnormal{(a)}]
Show that if~\(f\) is NOT uniformly continuous on~\(X\), then for some \(\varepsilon>0\) there are sequences \(\{p_n\}\), \(\{q_n\}\) in~\(X\) such that \(d_Y(f(p_n),f(q_n))>\varepsilon\) for each~\(n\) but \(d_X(p_n,q_n)\to 0\).
\item[\textnormal{(b)}]
Assume that~\(X\) is compact. Show, using (a), that~\(f\) is uniformly continuous on~\(X\).
\end{examsubparts}
\item
Let \(\{x_n\}\) be a sequence of points in \((a,b)\) and let \(\{c_n\}\) be a sequence of positive numbers such that \(\sum c_n\) converges. Define 
\[
f(x)=\sum_{n:x_n<x}c_n,\qquad a<x<b,
\]
 where the summation is understood as follows: sum over those indices~\(n\) for which \(x_n<x\). If there are no points \(x_n<x\), define the sum to be zero.
\begin{examsubparts}
\item[\textnormal{(a)}]
Show \(f(x-)=f(x)\) for each~\(x\) in \((a,b)\).
\item[\textnormal{(b)}]
Show \(f(x_n+)-f(x_n-)=c_n\), for each~\(n\).
\end{examsubparts}
\item
Suppose~\(f\) is continuous on \([0,\infty)\), differentiable on \((0,\infty)\), \(f(0)=0\) and \(f'\) is monotonically increasing. Define~\(g\) on \((0,\infty)\) by \(g(x)=f(x)/x\), \(x>0\). Prove~\(g\) is monotonically increasing.
\end{examproblems}
\end{document}
